\documentclass[a4paper,12pt]{article}
\usepackage{amsmath}
\usepackage{graphicx}
\usepackage{multicol}
\usepackage{geometry}
\usepackage{fancyhdr}
\usepackage{tikz}
\usepackage{eso-pic}
\usepackage[spanish]{babel}

\geometry{top=1.5cm, bottom=1.5cm, left=1.5cm, right=1.5cm}

\pagestyle{fancy}
\fancyhf{}
\rhead{Esc. 4-117 Ejército de los Andes}
\lhead{Termodinámica y Máquinas Térmicas}
\cfoot{\thepage}



\begin{document}

\begin{center}
    \Large \textbf{Evaluación de Termodinámica y Máquinas Térmicas} \\
    \normalsize Duración: 80 minutos \\
    Alumno/a: \underline{\hspace{6cm}} \hspace{0.5cm} Fecha: \underline{\hspace{2cm}}
\end{center}

\vspace{0.3cm}

\begin{multicols}{2}

\section*{Problemas}

\begin{enumerate}

\item \textbf{Trabajo de compresión en contexto real} \\
En una planta de producción de alimentos, un compresor toma $0{,}5\,\text{m}^3$ de aire a presión atmosférica y lo comprime hasta $0{,}02\,\text{m}^3$ a una presión final de $3\,\text{atm}$. Durante este proceso se realiza un trabajo de compresión de $8000\,\text{J}$. Calcule el trabajo de circulación total.

\vspace{0.3cm}

\item \textbf{Primer principio} \\
Un gas realiza un trabajo de expansión de $2562\,\text{kgm}$, manteniéndose constante su energía interna. Calcule el calor entregado al sistema, usando únicamente el primer principio de la termodinámica.

\vspace{0.3cm}

\item \textbf{Trabajo sobre gas } \\
Calcule el trabajo necesario para introducir $1{,}5\,\text{m}^3$ de gas en un reactor que opera a una presión constante de $6\,\text{atm}$. Explique brevemente la fórmula utilizada.

\vspace{0.3cm}

\item \textbf{Variación de energía interna} \\
Un gas es comprimido mediante un émbolo de área transversal $A = 0{,}01\,\text{m}^2$ que se desplaza $d = 0{,}3\,\text{m}$ bajo presión constante de $2\,\text{atm}$. Durante el proceso se transfieren $5\,\text{kcal}$ al ambiente. Calcule la variación de energía interna del gas.

\vspace{0.3cm}

\item \textbf{Trabajo de circulación} \\
Un compresor recibe $1\,\text{m}^3$ de aire a presión atmosférica y lo entrega comprimido a $0{,}025\,\text{m}^3$ y $3\,\text{atm}$, realizando un trabajo de compresión de $13500\,\text{kgm}$. Determine el trabajo de circulación total.

\end{enumerate}

\end{multicols}

\end{document}
